\documentclass[oneside,a4paper,14pt]{extarticle}
\usepackage[a4paper,top=20mm,bottom=20mm,left=30mm,right=10mm]{geometry}
\usepackage{amsmath, amssymb, amsfonts, amsthm}
\usepackage[T2A]{fontenc}
\usepackage[utf8]{inputenc}
\usepackage[russian]{babel}
\usepackage{textcomp}
\usepackage{indentfirst}
\usepackage{graphicx}
\usepackage{float}
\usepackage{wrapfig}
\usepackage{caption}
\usepackage{tocloft}
\renewcommand{\cftsecleader}{\cftdotfill{\cftdotsep}}
\usepackage[breaklinks]{hyperref}
\usepackage{titlesec}
\usepackage{fancyvrb}
\usepackage{xcolor}
\usepackage{minted}

\titleformat{\section}{\normalsize\bfseries}{\thesection}{1em}{}
\titleformat{\subsection}{\normalsize\bfseries}{\thesubsection}{1em}{}
\titleformat{\subsubsection}{\normalsize\bfseries}{\thesubsubsection}{1em}{}
\renewcommand\baselinestretch{1.33}\normalsize
\setlength{\parindent}{1.25cm}

\begin{document}
\newpage\thispagestyle{empty}
\begin{center}
	МИНИСТЕРСТВО НАУКИ И ВЫСШЕГО ОБРАЗОВАНИЯ\\
	РОССИЙСКОЙ ФЕДЕРАЦИИ\\
	ФЕДЕРАЛЬНОЕ ГОСУДАРСТВЕННОЕ БЮДЖЕТНОЕ\\
	ОБРАЗОВАТЕЛЬНОЕ УЧРЕЖДЕНИЕ ВЫСШЕГО ОБРАЗОВАНИЯ\\
	«ВЯТСКИЙ ГОСУДАРСТВЕННЫЙ УНИВЕРСИТЕТ»\\
	Институт математики и информационных систем\\
	Факультет автоматики и вычислительной техники\\
	Кафедра электронных вычислительных машин
\end{center}
\vspace{20mm}
\begin{center}
	Отчёт по лабораторной работе №4\\
	по дисциплине\\
	«Методы оптимизации»\\
\end{center}
\vfill
Выполнил студент гр. ИВТб-2301-05-00 \hfill
\makebox[8cm]{\hrulefill /Ковальский И.И./}
Проверил преподаватель\hfill
\makebox[8cm]{\hrulefill /Коржавина А.С./}
\vfill
\begin{center}
	Киров\\
	2026
\end{center}
\newpage
\setcounter{page}{1}
\tableofcontents
\newpage

\section*{Задание}
\addcontentsline{toc}{section}{Задание}
Реализовать один из методов решения задачи линейного программирования в соответствии с вариантом.\\

\textbf{Вариант:} Двойственный симплекс-метод (Dual Simplex) и метод ветвей и границ (Branch and Bound).\\

\textbf{Задание:}
\begin{enumerate}
	\item Изучить и реализовать заданный метод. Язык программирования выбрать самостоятельно.
	\item Продемонстрировать работу приложения на примерах. При демонстрации работы приложения выводить каждую итерацию метода (симплекс-таблицы) отдельно, также выводить решение (ответ). Предусмотреть все возможные варианты решения задачи, включая отсутствие решений и т.п.
\end{enumerate}

\newpage
\section*{Листинг кода}
\addcontentsline{toc}{section}{Листинг кода}

\begin{minted}[frame=lines,framesep=2mm,baselinestretch=1.2,fontsize=\footnotesize,breaklines]{python}

def knapsack_unbounded(weights, values, capacity):

    n = len(weights)
    dp = [0] * (capacity + 1)
    choice = [-1] * (capacity + 1)

    for cap in range(1, capacity + 1):
        best_val = dp[cap]
        best_idx = -1
        for i in range(n):
            w = weights[i]
            if w <= cap:
                cand = dp[cap - w] + values[i]
                if cand > best_val:
                    best_val = cand
                    best_idx = i
        dp[cap] = best_val
        choice[cap] = best_idx

    taken = []
    cap = capacity
    while cap > 0 and choice[cap] != -1:
        i = choice[cap]
        taken.append(i)
        cap -= weights[i]

    return dp[capacity], taken


def print_knapsack_table(weights, values, capacity):
    """Выводит таблицу DP для наглядности (как в транспортной задаче)."""
    n = len(weights)
    dp = [[0] * (capacity + 1) for _ in range(n + 1)]
    print(" № : вес : стоимость")
    for i, (w, v) in enumerate(zip(weights, values)):
        print(f"{i+1:2} :  {w:2}  :   {v:2}")
    print(f"Вместимость рюкзака: {capacity}\n")

    dp_unbounded = [0] * (capacity + 1)
    choices = [-1] * (capacity + 1)

    for cap in range(1, capacity + 1):
        best_val = 0
        best_i = -1
        for i in range(n):
            if weights[i] <= cap:
                cand = dp_unbounded[cap - weights[i]] + values[i]
                if cand > best_val:
                    best_val = cand
                    best_i = i
        dp_unbounded[cap] = best_val
        choices[cap] = best_i

    print(" виестимость | стоимость")
    for cap in range(capacity + 1):
        print(f"{cap:4} | {dp_unbounded[cap]:5}")

    taken = []
    cap = capacity
    while cap > 0 and choices[cap] != -1:
        i = choices[cap]
        taken.append(i)
        cap -= weights[i]

    print("\nВзятые предметы:", [i + 1 for i in taken])
    print("Общая стоимость:", dp_unbounded[capacity])
    print("Общий вес:", sum(weights[i] for i in taken))


def main_knapsack():
    weights = [1, 2, 7, 4]
    values  = [2, 5, 7, 9]
    capacity = 4

    print_knapsack_table(weights, values, capacity)

if __name__ == "__main__":
    main_knapsack()

\end{minted}

\newpage
\section*{Вывод}
\addcontentsline{toc}{section}{Вывод}

\end{document}